\documentclass[11pt]{article}

\usepackage{amsmath,amssymb,amsthm}
\usepackage{graphicx}
\usepackage{booktabs}
\usepackage[margin=1in]{geometry}
\usepackage[numbers]{natbib}
\usepackage{hyperref}

\newtheorem{theorem}{Theorem}
\newtheorem{lemma}[theorem]{Lemma}
\theoremstyle{definition}
\newtheorem{definition}[theorem]{Definition}

\title{{{title}}}
\author{Your Name \\ \small Your Institution \\ \small\texttt{you@example.org}}
\date{}

\begin{document}

\maketitle

\begin{abstract}
The problem, what was done, and what it showed, in a paragraph.
\end{abstract}

\section{Introduction}
\label{sec:introduction}

Why the problem matters, and what this paper adds to earlier work~\citep{knuth1984}.

\section{Method}
\label{sec:method}

\begin{definition}
  A term the rest of the paper uses.
\end{definition}

\begin{theorem}
  \label{thm:main}
  What is claimed.
\end{theorem}
\begin{proof}
  Why it holds.
\end{proof}

\section{Results}
\label{sec:results}

\begin{table}[ht]
  \centering
  \begin{tabular}{lrr}
    \toprule
    Method & Accuracy & Time (s) \\
    \midrule
    Baseline & 0.71 & 12.0 \\
    Ours     & 0.83 & 9.5 \\
    \bottomrule
  \end{tabular}
  \caption{What the table shows.}
  \label{tab:results}
\end{table}

Table~\ref{tab:results} compares them; Theorem~\ref{thm:main} explains why.

\section{Conclusion}
\label{sec:conclusion}

What was learnt, and what is left.

\bibliographystyle{plainnat}
\bibliography{references}

\end{document}
